\section{Related Work}

\subsection{Foundations}
InstructPix2Pix \cite{ref_ip2p} is the editor evaluated throughout this study. It fine-tunes a latent diffusion model \cite{ref_ldm} on a large corpus of synthetically generated editing triplets, produced by combining a language model with Prompt-to-Prompt-style attention injection \cite{ref_p2p}. Because supervision comes entirely from this generated data, the model inherits no explicit signal for preserving untouched content beyond what that generation process happened to keep stable, which is the mechanistic root of the drift problem this paper measures. CLIP \cite{ref_clip} supplies the embedding space the Drift Score measures change in, and the text-image similarity used to identify which region an instruction targets. Segment Anything \cite{ref_sam} supplies the region decomposition without which drift could only be measured at whole-image granularity, obscuring exactly the localized side effects this paper is concerned with.

\subsection{Consistency-Preserving Editing}
Prompt-to-Prompt \cite{ref_p2p}, Plug-and-Play \cite{ref_pnp}, and MasaCtrl \cite{ref_masactrl} are generation-time techniques for keeping an edited image consistent with its source, by manipulating cross-attention maps, injecting spatial features, or converting self-attention into a mutual query against the source image's own generation trajectory. All three aim to \textit{produce} a more consistent single edit. None of them measure how much unintended change remains afterward, and none address a chain of sequential edits. MasaCtrl's own motivating discussion states the semantic-drift problem in almost the same terms used here, but stops at fixing one edit rather than measuring or chaining it.

\subsection{Multi-Turn Editing Data}
MagicBrush \cite{ref_magicbrush} already contains manually annotated multi-turn editing sessions, which is independent validation that a four- or five-step chain is a realistic usage pattern rather than an artificial stress test. Its evaluation protocol, however, only checks whether each turn's target region changed the way it was supposed to. It never measures how much everything else changed in the process, which is exactly the gap this paper addresses.

\begin{table}[H]
\caption{Capability Comparison of Related Approaches}
\label{tab:capability}
\centering
\renewcommand{\arraystretch}{1.05}
\resizebox{\columnwidth}{!}{%
\begin{tabular}{@{}lcccc@{}}
\toprule
Approach & Unintended & Chain & Mitigation & Model \\
 & Change & Aware & Compare & Agnostic \\
\midrule
Prompt-to-Prompt \cite{ref_p2p}   & $\times$ & $\times$ & $\times$   & $\times$ \\
Plug-and-Play \cite{ref_pnp}      & $\times$ & $\times$ & $\times$   & $\times$ \\
MasaCtrl \cite{ref_masactrl}      & $\times$ & $\times$ & $\times$   & $\times$ \\
MagicBrush \cite{ref_magicbrush}  & Partial  & \checkmark & $\times$ & $\times$ \\
Proposed (Drift Score)            & \checkmark & \checkmark & \checkmark & \checkmark \\
\bottomrule
\end{tabular}%
}
\end{table}

No prior approach measures unintended regional change across a chain of edits or compares mitigation strategies on that specific axis (Table~\ref{tab:capability}); these dimensions are author-defined for this comparison rather than drawn from a standardized benchmark. The contribution of this paper is that combination: a measurement instrument built on top of the segmentation and embedding tools the field already uses, applied to the chain-based usage pattern the field already knows is realistic.
